\chapter{Proof Architecture}
\index{proof program!architecture|textbf}

\label{ch:proof-program}

\chapterstatus{proof contracts and proof strategies are presented with explicit status; no proposed physics proof package is represented as complete or reviewed.}


\section{Why a separate proof program is needed}

The computational hierarchy contains claims of several different kinds. Some
are elementary finite-dimensional identities, such as covariance under a
unitary coordinate change. Some require analytical hypotheses, such as a
spectral gap, compact parameter domain, resolvent bound, scale separation, or
asymptotic locality. Others are numerical or scientific statements that cannot
be established by proof alone.

A successful calculation may suggest a theorem but does not prove it. A formal
proof may establish an implication from assumptions while leaving the physical
adequacy of those assumptions unvalidated. The proof program therefore records

\begin{itemize}
  \item state spaces, domains, codomains, and group actions;
  \item assumptions needed by each result;
  \item analytical dependencies among results;
  \item mechanized theorem contracts and backend status;
  \item nonclaims and failure boundaries; and
  \item numerical or scientific evidence that remains separate from proof.
\end{itemize}

The proof workspace is subordinate to the versioned physical and numerical
specifications. It cannot repair an unresolved scientific convention by
choosing a mathematically convenient one.

\section{Proof-package graph}
\index{proof package!dependency graph}


The physical proof program is divided into six packages:

\begin{description}
  \item[\texttt{PRF-00}: Foundations.] State spaces, Bloch-fiber
    correspondence, and representation and reduction maps.
  \item[\texttt{PRF-05}: Mechanized operator lemmas.] Nine frozen
    finite-dimensional contracts targeted independently in Lean, Isabelle, and
    Rocq.
  \item[\texttt{PRF-10}: Operator alignment.] TB-anchored identification,
    gauge actions, and aligned pristine--doped subtraction.
  \item[\texttt{PRF-20}: Reduction.] Gauge-constrained locality,
    excluded-space reduction, atomistic-to-continuum consistency, and path
    commutativity.
  \item[\texttt{PRF-30}: Bounds.] Crossover, operator-to-observable, and
    continuum-fitting statements.
  \item[\texttt{PRF-40}: Compatibility.] Spectral/operator admissibility,
    model-class distance, and certified incompatibility.
\end{description}

\citationtodo{medium priority}{Add scholarly system references for Lean 4,
mathlib, Isabelle/HOL, and Rocq only after checking the sources and current
naming. Provisional keys: \texttt{demouraUllrich2021},
\texttt{mathlib2020}, \texttt{nipkowPaulsonWenzel2002}, and
\texttt{bertotCasteran2004}. These sources identify the systems; they cannot
substantiate repository-specific proof status.}

Their dependency structure is a partial order rather than a single pipeline:
foundations support every later package; mechanized elementary lemmas support
applicable alignment, reduction, and compatibility claims; alignment and
reduction support the proposed bounds; and compatibility also depends on
explicit model classes and admissible gauges.

The separate \texttt{PRF-90} package concerns agentic-workflow properties such
as authorization safety and replay determinism. It is intentionally outside the
physics proof chain. A workflow theorem would not establish a physical or
numerical theorem.

\section{Status vocabulary}
\index{proof!status vocabulary}


The proof registry distinguishes \texttt{proposed}, \texttt{in development},
\texttt{blocked}, \texttt{complete, unreviewed}, \texttt{reviewed},
\texttt{numerically checked}, and \texttt{retired}. Mathematical review,
numerical verification, scientific validation, and human publication acceptance
remain distinct.

A frozen mechanization contract fixes the current theorem target; it does not
mean that the theorem has been encoded or proved. A backend status of
\texttt{checked} means that the identified checker accepted the identified
encoding under the pinned toolchain. Cross-checking additionally requires all
three independent backends and semantic conformance review.

The derivations below explain selected contracts. They do not modify the owning
registry status.
